\chapter{Method}
\label{ch:method}

Chapter~\ref{ch:protocols} established that a Sapience request is broadcast in full, before
any fill, to every client connected to the relayer. This chapter sets out how that
disclosure is measured. It specifies a signalling model in which an informed forecaster
requests a quote from competing responders, defines the matched benchmark in which
responders may not condition on the request, states the two metrics, and closes by saying
what the design can and cannot establish. The experiment is a simulator, not a replica of
either venue: it is built to isolate one causal contrast, and its numbers describe that
contrast rather than the economics of any live market.

\section{The RFQ signalling model}
\label{sec:model}

A single binary event $Y \in \{0,1\}$ has public prior $\pi = P(Y=1)$. A contract pays one
unit if $Y=1$ and nothing otherwise.

\paragraph{The forecaster.} The forecaster observes a private signal $x$, drawn as
$x \mid Y=1 \sim \mathcal{N}(m, 1)$ and $x \mid Y=0 \sim \mathcal{N}(-m, 1)$, and forms the
posterior
\[
p_F \;=\; P(Y=1 \mid x) \;=\; \sigma\!\left(\log\tfrac{\pi}{1-\pi} + 2mx\right),
\]
where $\sigma$ is the logistic function. The signal strength $m$ is parameterised through
a signal-quality parameter $q_F = P(\operatorname{sign} x \text{ correct}) = \Phi(m)$.

One qualification belongs here rather than in a footnote. The forecaster acts on
$p_F > 0.5$, which coincides with the sign of $x$ only when $\pi = 0.5$; at other priors the
decision threshold shifts and the action's directional accuracy departs from $q_F$, while
the underlying signal strength $m$ is unchanged. With $q_F = 0.75$, for instance, the action
is correct $75.1\%$ of the time at $\pi = 0.5$ but $83.7\%$ at $\pi = 0.8$. $q_F$ should
therefore be read as signal quality, equal to directional accuracy at an even prior, and the
$\pi$ sweep of Section~\ref{sec:res-qf} varies the prior at fixed signal strength.

\paragraph{The action.} The forecaster does not reveal $p_F$. It submits a request
$R = (d, b)$, where $d = \mathbf{1}[p_F > 0.5]$ is the direction implied by $p_F$
and $b$ is a \emph{disclosed confidence band}: the confidence $|2p_F - 1|$ coarsened into
one of $B$ equal-width bands. The band count $B$ is the disclosure dial. At $B=1$ the
request carries direction alone. As $B$ grows, $R$ approaches full revelation of $p_F$.

The mapping onto the protocol is as follows. A Sapience broadcast carries the predicted
outcome on every leg and the position size \cite{sapience2026relayer}. Direction corresponds
to the predicted outcome; size stands in for the band, because a forecaster who is more
confident stakes more, so the disclosed size is a signal of confidence. A venue that
broadcasts an exact size therefore sits at the large-$B$ end of this range, and
Section~\ref{sec:cannot} returns to how tight that correspondence is.

\paragraph{The responders.} There are $n$ responders, indexed $j = 1 \dots n$. Responder $j$
observes a private binary signal $s_j$ that is correct with probability $q$. Signals are
exchangeable with dependence governed by $\rho$: with probability $\rho$ a responder's
signal is a copy of a single common draw, and otherwise it is drawn independently. $\rho$ is
therefore a mixture weight and not a correlation coefficient, and the signals are
conditionally independent given the event at $\rho=0$ rather than unconditionally
independent; the measured pairwise correlation there is $0.09$. Every signal is marginally correct with probability $q$ whatever $\rho$ is,
so $\rho$ varies how much \emph{independent} information the responders hold collectively
without varying how good any one of them is. That separation is what makes the $\rho$ sweep
interpretable.

\paragraph{Quoting.} Each responder quotes its own conditional posterior that $Y=1$,
truthfully. Without the request, responder $j$ quotes
\[
P(Y{=}1 \mid s_j) \;=\; \frac{\pi\,q^{s_j}(1-q)^{1-s_j}}
{\pi\,q^{s_j}(1-q)^{1-s_j} + (1-\pi)\,(1-q)^{s_j}q^{1-s_j}} .
\]
With the request, responder signals are independent of the forecaster's signal given $Y$, so
the action's likelihood multiplies in directly:
\[
P(Y{=}1 \mid R, s_j) \;=\;
\frac{P(R \mid Y{=}1)\,P(Y{=}1 \mid s_j)}
{P(R \mid Y{=}1)\,P(Y{=}1 \mid s_j) + P(R \mid Y{=}0)\,P(Y{=}0 \mid s_j)} .
\]
The action likelihoods $P(R \mid Y)$ are the probabilities that the forecaster's signal falls
in the region mapping to action $R = (d,b)$ under each state, and have no closed form because
the band boundaries are level sets of the posterior. They are computed once per parameter
cell by numerical quadrature of the two Gaussian densities over a fixed grid, and are a
function of the model's parameters alone rather than of any realised draw.

\paragraph{What responders are assumed to know, and what that costs.} Computing that
likelihood requires responder $j$ to know $q_F$, $B$ and $\pi$ exactly, so the model grants
them common knowledge of the forecaster's signal quality, of the disclosure granularity and
of the prior. This is a substantive assumption and it is not neutral: a responder that
misjudged $q_F$ would update too much or too little on the request, and the measured transfer
would fall. The assumption is made because the alternative, giving responders a prior over
$q_F$ and letting them learn it, is the cross-event learning the model deliberately excludes
and which Section~\ref{sec:cannot} lists as unmodelled. Its effect is therefore to place the
measurements at the upper end of what disclosure could be worth: they describe responders who
know exactly how informed their counterparty is.

\paragraph{Selection.} The forecaster then takes the best quote available: the lowest \textsc{yes} price when buying
\textsc{yes}, the highest when buying \textsc{no}.

This is best-of-$n$ posterior quoting and it is deliberately not a Bertrand equilibrium. It
is a behavioural rule, not a derived one, and two things follow that the reader should have
before any result. First, no responder shades, so none exercises market power at any $n$;
what varies with $n$ is only how many quotes the forecaster selects over. Second, and more
consequentially, quoting one's own posterior is zero-profit only unconditionally. The
responder that actually trades loses in expectation at every $n$ and in both arms, because
the model contains no uninformed order flow to offset losses to an informed counterparty,
and because under best-of-$n$ selection the winner is the responder whose signal was most
misleading in the forecaster's favour. A real market maker would price both effects by
widening its quote. This one cannot, and Section~\ref{sec:cannot} sets out what that costs
the interpretation of the result.

\section{The matched no-disclosure benchmark}
\label{sec:benchmark}

The experiment is a single causal contrast.

\begin{itemize}
  \item \textbf{Disclosure arm.} Responder $j$ observes $R$ and $s_j$, and quotes
    $P(Y=1 \mid R, s_j)$.
  \item \textbf{No-disclosure arm.} Responder $j$ observes $s_j$ only, and quotes
    $P(Y=1 \mid s_j)$.
\end{itemize}

Everything else is held identical: the same realised event, the same forecaster signal, the
same responder signals, the same $n$, $q$ and $\rho$, the same quoting rule and the same size
at risk. Both arms are run on the same drawn randomness, so the comparison is paired draw by
draw. The only difference between them is whether the responder's view carries the request.

\paragraph{What the benchmark is, and a second reading of $L$ that follows from it.} The
choice of benchmark determines what $L$ measures, so the consequence is stated here rather
than after the results. Because responders quote truthfully and the model contains no
uninformed order flow, the no-disclosure arm is a mechanism in which market makers lose
heavily and would in reality withdraw rather than quote. It is a clean counterfactual for
isolating the effect of disclosure, which is what it is for, and it is not a viable market.

Two readings of $L$ therefore survive and this report can exclude neither. On the first, $L$
is surplus transferred from the forecaster to the responders who quote. On the second, it is
the share of a structural loss that observing the request allows a responder to recover.
The arithmetic is identical under both; what differs is whether the baseline against which
the transfer is measured is one a responder would ever accept. Nothing in the experiment
distinguishes them, and Section~\ref{sec:limitations}(b) treats this as a limitation on what
the measurement can be said to establish.

\paragraph{The invariant, and how it is enforced.} Responders must observe the forecaster's
\emph{action} and never the private signal or the posterior. This is enforced structurally:
the quoting function receives a view object carrying the responder's own signal and, in the
disclosure arm, the direction and band. There is no field on that object through which $x$
or $p_F$ could travel, and the forecaster's state is a separate object never passed to a
responder.

\paragraph{Why a structural guarantee was not sufficient.} A first version of this model gave
the forecaster a binary signal and a deterministic threshold rule. The interface invariant
held: no responder function could reference the forecaster's signal. The invariant was
nonetheless vacuous, because with a binary signal and a threshold rule the action is a
lossless encoding of the signal. The direction \textsc{yes} occurred exactly when the
forecaster's signal was $1$, so a responder conditioning on the action had, in information
terms, conditioned on the signal itself; measured over $200{,}000$ draws, the action agreed
with the signal in every case, and the posterior took exactly two values. A model in which
the action encodes the signal without loss cannot measure disclosure, because there is no
disclosure decision left to vary. The defect was found by running the invariant check as
soon as the core loop first executed, rather than at the end, and it is reported here
because the fix is the reason the model has the shape it does: the continuous signal and the
banded action exist precisely so that $R$ is a strict coarsening of $p_F$. A parametrised
regression test asserts that the number of distinct actions is at most $2B$ while the number
of distinct posteriors is larger by more than an order of magnitude. That is a check on the
support of the action, not on how much information it carries: it would fail immediately if
the action space collapsed back onto the signal, which is the defect described here, and it
would pass under a subtler partial leakage. It bounds the cardinality and no more, and the
claim made for it should be no stronger.

\section{Parameters and sweep design}
\label{sec:sweep}

Six parameters are swept, one at a time, from the baseline $n=5$, $q=0.65$, $\rho=0$,
$q_F=0.75$, $B=8$, $\pi=0.5$, with $200{,}000$ draws per cell and the fixed seed
$20260907$. The baseline was fixed before any sweep was run, and the seed-robustness check
of Section~\ref{sec:res-seeds} repeats every sweep at the four further seeds $11$, $22$,
$33$ and $44$. Table~\ref{tab:params} lists them.

\begin{table}[htbp]
\centering
\begin{tabularx}{\textwidth}{p{0.10\textwidth}p{0.30\textwidth}p{0.20\textwidth}X}
\toprule
\textbf{Symbol} & \textbf{Meaning} & \textbf{Range swept} & \textbf{Why it is varied} \\
\midrule
$n$ & Number of competing responders & $1$ to $30$ & More responders means more parties observing the request and more competition for the business. Which dominates is the question \\
$q$ & Responder signal quality & $0.51$ to $0.95$ & A responder that is already well informed may learn little from the request \\
$\rho$ & Dependence among responder signals & $0$ to $1$ & Higher correlation means responders collectively hold less independent information \\
$B$ & Disclosed confidence bands & $1$ to $64$ & The disclosure dial: direction alone at $B=1$, approaching full revelation as $B$ grows \\
$q_F$ & Forecaster signal quality (directional accuracy at $\pi = 0.5$) & $0.55$ to $0.95$ & The headline magnitude depends on how well informed the forecaster is, so it must not be left at one value \\
$\pi$ & Public prior & $0.1$ to $0.9$ & An even prior is the most favourable case for an informed trader and should not be assumed \\
\bottomrule
\end{tabularx}
\caption{The swept parameters. $B$ is the one that varies the mechanism's design rather than
its environment, which is why Chapter~\ref{ch:evaluation} reads it as a mitigation.}
\label{tab:params}
\end{table}

\paragraph{What is not swept.} Two parameters named in the project's original design are
absent. Multi-leg requests, with a controlled redundancy parameter among legs, were not
implemented; the disclosure of joint belief argued in Section~\ref{sec:sapience} is
therefore argued and not measured. Order-flow concentration, the probability that the same
responder sees successive requests, was also not implemented. This matters more than it may
appear and Section~\ref{sec:cannot} treats it directly.

Two further experiments named in the original design were scoped out rather than left
unfinished, and are recorded here so that their later mention as future work is not a first
mention: a calibrated comparison against a pari-mutuel arm, cut because matching capital at
risk and transaction cost between an auction and a pool is a research question of its own;
and a specification-sensitivity experiment against a live judge, cut for time.

$q_F$ and $\pi$ were added to the sweep after a first pass reported a headline magnitude
that depended on both while holding both fixed. A result stated at one value of a parameter
that strongly governs it is not a result about the mechanism, and the two sweeps are
reported in Section~\ref{sec:res-qf}.

\section{Metrics}
\label{sec:metrics}

Two quantities are reported together.

\paragraph{The transfer.} Buying \textsc{yes} at price $P$ the forecaster receives $Y - P$;
buying \textsc{no} it pays $1 - P_{\textsc{yes}}$ and receives $(1-Y)$ less that price, so in
both directions it gains when the contract settles its way. The economic measure is
\[
L \;=\; S_{\text{no-disclosure}} \;-\; S_{\text{disclosure}},
\]
where $S$ is the forecaster's expected surplus, that is the realised payoff net of the price
paid. Chapter~\ref{ch:results} also reports $L$ as a proportion,
$L / S_{\text{no-disclosure}}$, the share of the forecaster's benchmark surplus that does not
survive disclosure. $L$ is a measured quantity with a deliberately neutral name. It is not called the MEV,
in the code, in the figures or in the prose, because how much of it meets the definition in
Chapter~\ref{ch:introduction} is the question Chapter~\ref{ch:evaluation} asks, and that
question must not be settled by a variable name.

Because both arms run on identical randomness, $L$ is a paired difference and its standard
error is that of the per-draw difference rather than the sum of the two arms' errors. At
$200{,}000$ draws the $95\%$ interval on $L$ is $\pm 0.0007$ at the baseline, and
Chapter~\ref{ch:results} draws it on every figure.

That interval answers whether $L$ differs from zero in a given cell. It is the wrong
interval for asking whether $L$ differs between two cells. Every cell in a sweep is run from
the same seed and consumes the random stream in the same order, so cells share the realised
event and the forecaster's signal, and a between-cell difference is itself paired. Its
interval is typically an order of magnitude tighter: comparing $n=15$ with $n=30$, for
instance, the correct interval is $\pm 2\times10^{-5}$ rather than $\pm 7\times10^{-4}$;
that comparison was computed directly for this sentence, since the recorded step intervals
are between adjacent cells.
The sweeps therefore record, for every step, both the change in $L$ and the paired interval
on that change, and Chapter~\ref{ch:results} judges resolution against the latter. Using a
single cell's interval for this purpose understates the experiment's resolving power, which
is a conservative error but an error nonetheless.

\paragraph{The information.} The informational measure is how much a responder learns from
the request,
\[
I \;=\; \mathbb{E}_R\!\left[\, D_{KL}\!\left(P(Y \mid R, s_j) \,\big\|\, P(Y \mid s_j)\right) \right],
\]
reported in nats, and averaged over responders as well as over draws. Each per-draw
divergence is exact, because the simulator has the true posteriors, but the reported figure
is a Monte Carlo mean like any other and its interval is recorded alongside it.

Reporting $I$ alongside $L$ is intended to separate how much is learned from how much
surplus moves. Section~\ref{sec:res-info} shows that $I$ is invariant by construction in two
of the six sweeps, so it contributes as a check in those and as a measurement only in the
other four.

The quoted spread, defined on each draw as the difference between the highest and lowest
of the $n$ quotes and averaged over draws, both arms' surpluses and the responder surplus
are also recorded. Because
the model is zero-sum between the two parties, responder surplus is the negative of
forecaster surplus and no crossover between them is possible; that question, raised in the
project's original design, is therefore answered trivially rather than empirically.

\section{Implementation and reproducibility}
\label{sec:implementation}

The simulator is about 450 lines of Python using NumPy and SciPy, excluding tests. Every run is
seeded and deterministic. Sweeps write CSV files to \texttt{data/raw}, and a separate script
turns those files into the figures reproduced in Chapter~\ref{ch:results}; no figure in this
report is edited by hand, so each is regenerable with a single command. Eleven regression tests
cover the model, including two posterior calculations checked by hand against the arithmetic
in Section~\ref{sec:model}, a reconstruction of the degenerate binary-signal configuration
described above, a check that the action likelihoods form a distribution, the
parametrised coarsening assertion described above, and a test that quotes in the
no-disclosure arm do not move when the action changes.

Because every reported number comes from one realisation, the sweeps were also re-run in
full at five seeds; Section~\ref{sec:res-seeds} reports the result. The repository records the sweep parameters and the
seed in every data file, and the date and configuration of the run alongside them.

Three numerical guards are applied and are noted here because they are not otherwise
visible: the action-likelihood table is floored at $10^{-15}$ so that an unreachable band
cannot produce a division by zero, the divergence calculation clips its arguments away from
zero and one, and the band index is clipped to its valid range for the case where a
posterior saturates to exactly one in floating point. All three are guards rather than
adjustments and none changes a reported quantity, but the implementation is described here
as complete and they belong in that description.

\section{What this design can and cannot establish}
\label{sec:cannot}

This section is deliberately placed before the results rather than after them.

\paragraph{The sign of the effect is expected, and is not the finding.} That allowing a
market maker to condition on an informed order moves surplus from the informed party to the
market maker is the standard adverse-selection result. Glosten and Milgrom show that a
specialist facing a mixture of informed and uninformed order flow sets a spread reflecting
the probability that the counterparty is informed \cite{glosten1985bidask}, and Kyle shows
that an informed trader's advantage is bounded by how much the market maker can infer from
the order flow it observes \cite{kyle1985continuous}. A simulator that failed to reproduce
$L > 0$ would be evidence that the simulator was wrong, not evidence about disclosure.
Chapter~\ref{ch:results} therefore treats the sign as a validity check on the harness and
reports the \emph{shape} of the effect as the finding: its direction and monotonicity in
$n$, its behaviour in $q$, its response to $\rho$, and the relationship between $I$ and $L$.

\paragraph{The numbers are internal to the model.} $L$ is expressed in units of the binary
contract under a specific signal structure, prior and competition rule. It is not an estimate
of value transferred on any live venue, and no such estimate is offered anywhere in this
report. What transfers across is the comparative structure: which direction the effect moves
as a parameter changes, and where it is largest.

\paragraph{The accumulation limb is argued, not demonstrated.} Chapter~\ref{ch:protocols}
observed that an address in every broadcast makes a per-address history cheap to
accumulate. The simulator does not model this. Responders carry no state across events, so
even had the concentration parameter been implemented it would have measured repeated
\emph{exposure} rather than accumulated inference. Whether accumulation converts into an
advantage that competitive quoting does not compensate is therefore a proposition this
report argues from protocol structure, and it is one of the two places where the evidence is
weakest.

\paragraph{The sizing decision is abstracted away.} Traded size is held constant across all
cells. An earlier configuration scaled size with the disclosed band, which confounded the
disclosure dial: raising $B$ then both revealed more and changed how much the forecaster
staked, and the no-disclosure surplus moved with a parameter that should not have touched it.
Holding size fixed makes $B$ vary disclosure and nothing else, at the cost of removing the
forecaster's choice of how much to stake. A forecaster who reduced size to reduce disclosure
is therefore outside the model, and the correspondence between $B$ and a real broadcast size
is weaker for it.

\paragraph{Responders quote truthfully, bear no risk, and cannot widen.} Real market makers
hold inventory, face capital constraints, price risk as well as information, and shade their
quotes against the possibility that winning is itself bad news. The model's responders do
none of this, so the surplus attributed here to information would in a real venue be mixed
with compensation for risk borne. That is the distinction the report's definition turns on,
and the model resolves it by construction rather than by measurement.

\paragraph{There is no uninformed order flow, so the benchmark is not a viable mechanism.}
Every request in this model comes from the informed forecaster. In Glosten and Milgrom a
market maker survives because gains from uninformed traders offset losses to informed ones
\cite{glosten1985bidask}; with that population removed, responders lose in every
configuration and would exit rather than quote. Two consequences follow. The magnitudes are
upper bounds on the forecaster's position rather than estimates of any real one. And $L$
admits a second reading that this report cannot exclude: measured against a benchmark no
responder would join, it is as defensible to call it the share of a structural loss that
disclosure recovers as to call it surplus transferred. Chapter~\ref{ch:evaluation} returns
to this.

\paragraph{One event, one forecaster, one round.} There is no repeated play, no reputation,
no competition between forecasters, and no strategic choice of whether to trade at all. A
forecaster who could decline to submit a request when the expected disclosure cost exceeded
the expected gain would face a different problem from the one modelled here.

What the design does establish is a clean causal contrast under a matched benchmark: holding
competition, signal quality and dependence identical, with no risk borne in either arm, it
isolates the effect of
letting responders condition on the request, and it measures how that effect varies with the
environment and with the granularity of what is disclosed. Chapter~\ref{ch:results} reports
those measurements and Chapter~\ref{ch:evaluation} asks what they are worth.
